\chapter{Introduction}
\label{ch:introduction}

\section{Motivation}
\label{sec:motivation}

Project management platforms such as Trello, Jira, Asana, GitHub Projects, and GitLab are the default tools used by software teams today to plan, assign, and track work (\Cref{sec:traditional-pm}). These platforms are mature and effective at their core task, but they share an architectural property that is easy to overlook: every project record, every status change, and every audit trail they produce is stored and controlled entirely by a single provider. If a team wants to verify \emph{independently} that a task's history has not been altered, there is no mechanism to do so — they can only trust that the platform has not modified it.

This matters in any collaborative setting where trust between participants cannot be assumed, or where an independently verifiable record of who did what, and when, is itself a requirement rather than a convenience. Multi-organization software projects, outsourced development work, and academic or regulatory contexts where accountability must be demonstrable are all cases where a centrally-controlled audit log is a weaker guarantee than participants may need.

Blockchain technology, and permissioned blockchain frameworks such as Hyperledger Fabric in particular (\Cref{sec:fabric}), were designed to solve exactly this class of problem: providing a shared, tamper-evident record of state changes among a known set of participants, without requiring any single participant to be fully trusted. \textcite{adel2023} demonstrate that this approach is practical for collaborative data management by applying it to the construction industry (\Cref{sec:adel}), combining Hyperledger Fabric for on-chain metadata with the InterPlanetary File System (IPFS) for off-chain document storage. However, their system targets a construction-specific data model that does not map onto how software teams organize their work.

This thesis addresses that gap: it applies the validated Fabric-plus-IPFS architectural pattern to the specific domain of software project management, producing a system that offers the functional surface software teams already expect — projects, tasks, priorities, assignments, status boards, file attachments — while replacing the centrally-controlled audit trail of existing tools with an immutable, independently verifiable one.

\section{Content of the Thesis}
\label{sec:content}

\Cref{ch:background} introduces the concepts needed to understand the rest of the document: the structure and requirements of software project management (\Cref{sec:pm}), the limitations of existing centralized platforms (\Cref{sec:traditional-pm}), and the technical foundations — blockchain, smart contracts, Hyperledger Fabric, and IPFS — that the proposed system builds on.

\Cref{ch:related} discusses \textcite{adel2023} in depth, the closest prior work to this thesis, and uses it to precisely position this thesis' contribution relative to both the blockchain-based data management literature and the traditional project management tools discussed in \Cref{ch:background}.

\Cref{ch:method} describes the system itself: its layered architecture (\Cref{sec:architecture}), the chaincode that governs project and task state on the Fabric ledger (\Cref{sec:chaincode-design}), the integration with IPFS for file attachments (\Cref{sec:ipfs-integration}), the backend API that exposes this functionality over HTTP (\Cref{sec:backend}), and the frontend interface used to interact with it (\Cref{sec:frontend}). This chapter also discusses a concrete implementation issue encountered during development — a chaincode determinism bug caused by including timestamps generated independently by each peer — as an illustration of a constraint that is easy to overlook when building Fabric-based systems.

\Cref{ch:results} evaluates the resulting prototype against the requirements set out in \Cref{sec:pm}, demonstrating the end-to-end workflow — from project creation through task assignment, status transitions, and file attachment — and examining the immutable audit trail the system produces as its central contribution.

Finally, \Cref{ch:conclusion} summarizes what was achieved, discusses the limitations of the current prototype, and outlines directions for future work.
